% TOP_PROBLEM: {"id": 2309007, "title": "Research Problems in Function Theory — Problem 9.7", "category_id": 8, "category": {"id": 8, "name": "analysis", "display_name": "Analysis", "description": "Limits, continuity, calculus, and function theory.", "slug": "analysis", "order_index": 8, "created_at": "2026-07-31T15:26:25.671Z"}, "status": "open"}
% FIRST_POSTED: null
% ENTRY_KIND: statement_only
% Imported from: batch07.tex; UnsolvedMath ID 2309007
\documentclass[11pt]{book}
\usepackage{fontspec}
\setmainfont{Latin Modern Roman}
\setsansfont{Latin Modern Sans}
\usepackage{amsmath,amssymb,mathtools}
\usepackage{unicode-math}
\setmathfont{Latin Modern Math}
\usepackage{xurl}
\usepackage[hidelinks]{hyperref}
\newcommand{\sref}[2]{\hyperlink{source:#1:#2}{[S#2]}}
\newcommand{\cataloguescope}[1]{\paragraph{Mathematical question.}#1}
\begin{document}
% UnsolvedMath ID 2309007; source catalogue code AMR-022-9007
\setcounter{section}{2309006}
\section{Research Problems in Function Theory — Problem 9.7}\label{2309007}
{\small\sffamily UnsolvedMath ID 2309007 · Analysis}
\subsection{Definitions and mathematical statement}

\paragraph{Shared notation.}$\mathbb N=\{1,2,\ldots\}$, and $\mathbb Z,\mathbb Q,\mathbb R,\mathbb C$ denote the integers, rationals, reals and complex numbers. Any other conventions are specified locally.


For a closed $E\subset\mathbb C$, let $A(E)=\{g\in C(E):g\text{ is analytic on }E^{\circ}\}$. Call $E$ weak Arakelian if for every $g\in A(E)$ there is an entire function $F$ such that, for every sequence $(z_j)\subset E$, \[|F(z_j)|\to\infty\quad\Longleftrightarrow\quad |g(z_j)|\to\infty.\] The roles of the prescribed function and the entire function are separated here by the names $g$ and $F$.
\paragraph{Statement recorded in the supplied source.}
Give necessary and sufficient geometric conditions for a closed planar set $E$ to be weak Arakelian. \sref{2309007}{2}
\subsection{Short English statement}
Characterize closed sets on which each continuous interior-analytic function has the same diverging-value sequences as some entire function.
\subsection{Sources}
\begin{description}
\item[\hypertarget{source:2309007:1}{S1}] ULAM.AI and the UnsolvedMath Contributors, \emph{UnsolvedMath: A Curated Collection of Open Mathematics Problems}, record 2309007 (AMR-022-9007). CC BY 4.0. \url{https://huggingface.co/datasets/ulamai/UnsolvedMath}.
\item[\hypertarget{source:2309007:2}{S2}] Walter K. Hayman and Edward F. Lingham, \emph{Research Problems in Function Theory} (2018), arXiv:1809.07200. Problem 9.7, complete original question, all following updates, and the relevant chapter definitions. \url{https://arxiv.org/abs/1809.07200}.
\end{description}
\label{end:2309007}
\end{document}
